In one period \(T\), a wave crest travels one wavelength \(\lambda\). Since \(f=1/T\), the wavespeed relation is

\[c=\frac{\lambda}{T}=\lambda f,\qquad f=\frac{c}{\lambda},\qquad \lambda=\frac{c}{f}.\]

The supplied speeds in SI units are \(c_{\mathrm{light}}=3.00\times10^8\ \mathrm{m/s}\) and \(c_{\mathrm{sound}}=340\ \mathrm{m/s}\). These values are used throughout; unit prefixes are converted before substitution.

For A.2.1, the visible-light wavelength is \(\lambda=510\ \mathrm{nm}=510\times10^{-9}\ \mathrm{m}\). Therefore,

\[\boxed{f=\frac{3.00\times10^8}{510\times10^{-9}}\ \mathrm{Hz}=5.88\times10^{14}\ \mathrm{Hz}=588\ \mathrm{THz}}\qquad\text{(A.2.1)}\]

For A.2.2, the AM wavelength is \(\lambda=550\ \mathrm{m}\). Using the same relation gives

\[\boxed{f=\frac{3.00\times10^8}{550}\ \mathrm{Hz}=5.45\times10^5\ \mathrm{Hz}=545\ \mathrm{kHz}}\qquad\text{(A.2.2)}\]

For A.2.3, the infrared wavelength is \(\lambda=10\ \mu\mathrm{m}=10\times10^{-6}\ \mathrm{m}\). Its frequency is

\[\boxed{f=\frac{3.00\times10^8}{10\times10^{-6}}\ \mathrm{Hz}=3.00\times10^{13}\ \mathrm{Hz}=30.0\ \mathrm{THz}}\qquad\text{(A.2.3)}\]

The next item supplies frequency instead of wavelength. For A.2.4, \(f=89\ \mathrm{MHz}=89\times10^6\ \mathrm{Hz}\), so

\[\boxed{\lambda=\frac{3.00\times10^8}{89\times10^6}\ \mathrm{m}=3.37\ \mathrm{m}}\qquad\text{(A.2.4)}\]

For A.2.5, the observed delay between the flash and thunder is \(\Delta t=36-32=4\ \mathrm{s}\). Taking the light and sound to originate at the same instant and travel the same distance \(d\), their arrival-time difference is

\[\Delta t=\frac{d}{c_{\mathrm{sound}}}-\frac{d}{c_{\mathrm{light}}}.\]

Solving this equation for distance accounts for both travel times and gives

\[d=\frac{\Delta t}{1/c_{\mathrm{sound}}-1/c_{\mathrm{light}}}
=\frac{4}{1/340-1/(3.00\times10^8)}\ \mathrm{m}.\]

Using \(1\ \mathrm{km}=1000\ \mathrm{m}\) and the exact conversion \(1\ \mathrm{mi}=1609.344\ \mathrm{m}\), the required results are

\[\boxed{d=1.360\ \mathrm{km}=0.845\ \mathrm{mi}}\qquad\text{(A.2.5)}\]

Because light travels much faster than sound, the usual approximation \(d\simeq c_{\mathrm{sound}}\Delta t=340(4)=1360\ \mathrm{m}\) gives the same reported answers. The finite-light correction is only \(1.541\ \mathrm{mm}\).

For A.2.6, the supplied mobile-signal frequency is \(f\approx29\ \mathrm{GHz}=29\times10^9\ \mathrm{Hz}\). Its wavelength is

\[\boxed{\lambda=\frac{3.00\times10^8}{29\times10^9}\ \mathrm{m}=10.34\ \mathrm{mm}}\qquad\text{(A.2.6)}\]

The microwave oven in A.2.7 operates at \(f=2.5\ \mathrm{GHz}=2.5\times10^9\ \mathrm{Hz}\). Applying the same wavespeed relation answers part (a),

\[\boxed{\lambda=\frac{3.00\times10^8}{2.5\times10^9}\ \mathrm{m}=0.1200\ \mathrm{m}=120.0\ \mathrm{mm}}\qquad\text{(A.2.7.a)}\]

For part (b), the problem specifies \(D\leq\lambda/2\). Consequently, the largest diameter allowed by that stated criterion is

\[\boxed{D_{\max}=\frac{\lambda}{2}=\frac{0.1200}{2}\ \mathrm{m}=60.0\ \mathrm{mm}}\qquad\text{(A.2.7.b)}\]

For A.2.7.c, we measured the mesh on our Oster small office microwave and found a hole diameter of roughly \(1.2\ \mathrm{mm}\) and a center-to-center pitch of roughly \(1.6\ \mathrm{mm}\). Using the measured diameter gives

\[\boxed{D_{\mathrm{mesh}}\approx1.2\ \mathrm{mm},\qquad
\frac{D_{\mathrm{mesh}}}{\lambda}\approx0.01000,\qquad
\frac{D_{\max}}{D_{\mathrm{mesh}}}\approx50.0}\qquad\text{(A.2.7.c)}\]

Thus, the measured hole is approximately one hundredth of the microwave wavelength and fifty times smaller than the classroom upper bound. Its opening is nevertheless much larger than a visible-light wavelength, so food remains visible through the perforated screen. The \(60.0\ \mathrm{mm}\) result answers the supplied \(\lambda/2\) criterion; it is not a microwave-door design specification.

For the graduate calculation in A.2.8.a, use the stipulated length \(D=2\ \mathrm{pm}=2\times10^{-12}\ \mathrm{m}\) as the electromagnetic wavelength. Then

\[\boxed{f=\frac{c_{\mathrm{light}}}{D}=\frac{3.00\times10^8}{2\times10^{-12}}\ \mathrm{Hz}=1.500\times10^{20}\ \mathrm{Hz}=150.0\ \mathrm{EHz}}\qquad\text{(A.2.8.a)}\]

This calculation uses the supplied length scale without establishing a physical electron diameter. The resulting wavelength is \(0.00200\ \mathrm{nm}\), below the \(0.01\ \mathrm{nm}\) gamma-ray range described by NASA. Therefore, the wavelength-based classification requested in A.2.8.b is

\[\boxed{\text{Gamma ray}}\qquad\text{(A.2.8.b)}\]

\clearpage
\matlabcaption{Script used to reproduce the numerical results in A.2.}

\lstinputlisting[style=hwmatlab]{matlab/A2.m}
